\documentclass[12pt]{article}
\usepackage[T1]{fontenc}
\usepackage[utf8]{inputenc}
\usepackage{lmodern}
\usepackage{textcomp}
\usepackage{enumitem}
\usepackage{geometry}
\usepackage{graphicx}
\usepackage{amsmath, amssymb}
\geometry{margin=1in}
\begin{document}
\begin{center}
Political Science - Undergraduate Level Assessment 19 \
Total Marks: 40 \
Answer all questions.
\end{center}

\section*{Multiple Choice (10 marks)}
\begin{enumerate}[label=\arabic*.]
\item The realm of policy decisions concerned primarily with relations between the United States and the rest of the world is known as
\begin{enumerate}[label=(\alph*)]
    \item terrorism policy.
    \item economic policy.
    \item foreign policy.
    \item international policy.
\end{enumerate}
\item All of the following are constitutionally mandated presidential powers EXCEPT the power to
\begin{enumerate}[label=(\alph*)]
    \item negotiate treaties
    \item veto legislation
    \item grant reprieves and pardons
    \item declare war
\end{enumerate}
\item Which of the following does NOT appear in the Constitution?
\begin{enumerate}[label=(\alph*)]
    \item The electoral college
    \item Political parties
    \item Separation of powers
    \item The term length for members of Congress
\end{enumerate}
\item Which of the following are possible constraints on US foreign policy decision making?
\begin{enumerate}[label=(\alph*)]
    \item Foreign policies of other states
    \item International law
    \item Intergovernmental organizations
    \item All of the above
\end{enumerate}
\item What were the implications of the Cold War for American exceptionalism?
\begin{enumerate}[label=(\alph*)]
    \item It ended the influence of American exceptionalism entirely
    \item Exceptionalism was enhanced by America's status as the 'leader of the free world'
    \item The extension of American power globally challenged core assumptions of exceptionalism
    \item Both b and c
\end{enumerate}
\end{enumerate}

\section*{True / False (10 marks)}
\textit{Answer True or False}
\begin{enumerate}[label=\arabic*.]
\setcounter{enumi}{5}
\item True or False: The term 'iron triangle' refers to the stable relationship between federal bureaucracy, congressional committees, and lobbyists.
\item True or False: Since the 1980 s, the Republican Party has been dominated by labor unions, shaping its political agenda significantly.
\item True or False: The Voting Rights Act of 1965 mandated that all voters must pass literacy tests before being allowed to vote.
\item True or False: Strict constructionism argues that the Bill of Rights only limits the powers of the federal government and not the states.
\item True or False: The Voting Rights Act of 1965 suspended the use of literacy tests at voting centers to enhance voter access.
\end{enumerate}

\section*{Long Form Response (20 marks)}
\textit{Answer each question with a clear and complete explanation.}
\begin{enumerate}[label=\arabic*.]
\setcounter{enumi}{10}
\item Discuss the argument made by James Madison in The Federalist No. 10 regarding the impact of factionalism and how a republican government can mitigate its negative effects.
\item In what ways did the geopolitical dynamics of the Cold War influence U.S. perceptions and policies towards developments in the Third World? Discuss with specific examples.
\end{enumerate}

\vfill
\noindent\textit{End of Paper}
\end{document}